% Results section.
%
% EVERY NUMBER IN THIS FILE IS A MACRO OR A GENERATED TABLE. Nothing is typed
% literally. The macros are defined in paper/generated/results_macros.tex and
% the tables are emitted by:
%
%   python training/holdout_eval.py --manifest <manifest> \
%       --writeup training/run_records/last_run_v2b_*.json \
%                 training/run_records/last_run_v2c_*.json \
%                 training/run_records/last_run_v2d_*.json \
%       --trainable 3697923 10795779 7246851 \
%       --confusion <matrix.json> --tex-out paper/generated
%
% If a figure here disagrees with a run record, the run record is right and
% this file was edited by hand, which it must not be.

\section{Results}
\label{sec:results}

\subsection{Reported model}

We report the configuration that trains the top four encoder layers at a peak
learning rate of $1\times10^{-5}$. Early stopping on a held-out stopping set
ended the run before its step ceiling and selected step \ResultSelectedStep;
the ceiling and the stopping step are recorded in
\texttt{training/run\_records/}. Table~\ref{tab:results} gives its
performance on the reporting set, which no stopping decision has observed.

% GENERATED FILE - DO NOT EDIT BY HAND.
% Written by training/evaluate.py --writeup in ONE evaluation pass.
% Editing this file to change a reported number is fabrication;
% re-run the emitter instead.
%
% reported_model: model_v2d_freeze8_lr1e-5
% run_fingerprint: 3fdcb8c5d5f8f137
% manifest: dataset/processed/eval_manifest_phrase_v2.json
% strategy: phrase
% split_seed: 42
% reporting_rows: 17942
% reporting_phrases: 9
% reporting_groups: 5
% max_length: 96
% git_commit: ed02093
% generated_at: 2026-09-08T04:59:19.668469+00:00
%
\begin{table}[t]
\centering
\caption{Triage performance of the reported model on the frozen phrase holdout. The 17,942 rows are frame permutations of 9 distinct phrases in 5 phrase groups; the support column is therefore not an independent sample size, and the final column gives the count that governs the granularity of each recall figure.}
\label{tab:results}
\begin{tabular}{lrrrrr}
\toprule
Class & Precision & Recall & F1 & Rows & Distinct sentences \\
\midrule
CRITICAL & 0.6633 & 0.8504 & 0.7453 & 8,962 & 4 \\
URGENT & 0.7426 & 0.4963 & 0.5950 & 7,793 & 4 \\
ROUTINE & 0.9549 & 1.0000 & 0.9770 & 1,187 & 1 \\
\midrule
Macro avg & --- & --- & 0.7724 & 17,942 & 9 \\
\bottomrule
\end{tabular}
\\[2pt]{\scriptsize Run fingerprint \texttt{3fdcb8c5d5f8f137}}
\end{table}


Macro F1 is \ResultMacroFOne. CRITICAL recall is \ResultCriticalRecall\ at a
precision of \ResultCriticalPrecision; URGENT recall is \ResultUrgentRecall.
The model does not meet our acceptance gate: it fails the CRITICAL recall
condition while satisfying both non-degeneracy conditions
(Section~\ref{sec:gate}).

The support column should not be read as a sample size. The
\ResultEvalRows\ reporting rows are frame permutations of
\ResultEvalPhrases\ distinct phrases in \ResultEvalGroups\ phrase groups, and
the rightmost column of Table~\ref{tab:results} gives the count that actually
governs each recall figure. We return to this in
Section~\ref{sec:limitations}, because it bounds what any of these numbers can
mean.

\paragraph{Comparison against human practice: \PENDING{no nurse study has
run}.} These figures are agreement with a label set, not with practice. What a
nurse assigns given the same description has not been measured, so the model is
unanchored against the thing it would assist. The study design --- case
selection at the concept level, blinding, and per-class analysis with
over-triage and under-triage reported separately rather than averaged --- is
written and unexecuted.

\subsection{Where the error is}

% GENERATED FILE - DO NOT EDIT BY HAND.
%
\begin{table}[t]
\centering
\caption{Confusion matrix for v2d\_freeze8\_lr1e-5 on the reporting set. ROUTINE is separated perfectly; the residual error is entirely on the CRITICAL/URGENT boundary, and this was true of every configuration trained.}}
\label{tab:confusion}
\begin{tabular}{lrrr}
\toprule
Truth $\downarrow$ / Pred $\rightarrow$ & CRITICAL & URGENT & ROUTINE \\
\midrule
CRITICAL & 7,621 & 1,341 & 0 \\
URGENT & 3,869 & 3,868 & 56 \\
ROUTINE & 0 & 0 & 1,187 \\
\bottomrule
\end{tabular}
\end{table}


Table~\ref{tab:confusion} locates the residual error precisely. ROUTINE is
separated without error in both directions. Every misclassification the model
makes lies on the CRITICAL/URGENT boundary, in both directions: CRITICAL
presentations read as URGENT, and URGENT presentations read as CRITICAL. This
was true of every configuration we trained, not only the reported one.

\subsection{Configuration search}

% GENERATED FILE - DO NOT EDIT BY HAND.
% Written by training/evaluate.py --writeup from saved run records.
%
\begin{table}[t]
\centering
\caption{Three configurations on the same frozen phrase holdout, same split seed, same reporting set. Trainable parameters exclude the frozen 250{,}002$\times$384 embedding table, which is 96{,}199{,}296 of the model's 117{,}641{,}859 parameters in every row.}
\label{tab:sweep}
\begin{tabular}{lrrrrrr}
\toprule
Run & Trainable & Sel.\ step & Macro F1 & CRIT P/R & URG P/R & Gate \\
\midrule
v2b\_freeze10\_lr1e-5 & 3,697,923 & 1200 & 0.7660 & 0.7951/0.9007 & 0.8636/0.6315 & FAIL \\
v2c\_freeze6\_lr1.5e-5 & 10,795,779 & 450 & 0.5545 & 0.5337/0.9749 & 0.2321/0.0083 & FAIL \\
v2d\_freeze8\_lr1e-5 & 7,246,851 & 900 & 0.7724 & 0.6633/0.8504 & 0.7426/0.4963 & FAIL \\
\bottomrule
\end{tabular}
\end{table}


Table~\ref{tab:sweep} reports three configurations evaluated on the same
frozen holdout, the same split seed and the same reporting set, differing only
in how much of the encoder was left trainable and in the peak learning rate.
The table's caption gives the frozen embedding count, which is identical in
every row, so the trainable column is the quantity that varies.

Two observations, and we draw a different conclusion from each.

The first is that trainable capacity governs \emph{when} the loss curve turns.
Holding the learning rate fixed, the middle configuration's loss inflected
several hundred steps earlier than the smallest one's, and the largest
configuration's earlier still. The per-evaluation stopping-loss curves that
show this are committed in \texttt{training/run\_records/}; we do not quote
inflection steps here because we have no principled estimator for them, only
a reading of the curves. The effect is what makes the smallest configuration
expensive rather than wrong.

The second is that capacity does not change \emph{what} the model gets wrong.
Across a threefold range of trainable parameters and a $1.5\times$ range of
learning rate, every configuration separated ROUTINE perfectly and none
resolved the CRITICAL/URGENT boundary. We therefore stopped searching, and we
report the search as inconclusive on that boundary rather than presenting the
best row as a tuning result. Section~\ref{sec:limitations} gives the
measurement-granularity argument for why continuing would not have been
informative.

\subsection{A single-metric gate certified a model that had abandoned a class}

% GENERATED FILE - DO NOT EDIT BY HAND.
% Written by training/evaluate.py --writeup in ONE evaluation pass.
% Editing this file to change a reported number is fabrication;
% re-run the emitter instead.
%
% reported_model: model_v2d_freeze8_lr1e-5
% run_fingerprint: 7c40d09ccb91f0d1
% source_sha256: 2ea0d2d495617e6c7385faa3afce1c7ccfcedfa2bc438d75620ba8e98a40b394
% eval_sha256: b8cddf0dfd1646856ae44637dd724c542b4203fa159e150d01e42efc39723af5
% manifest: dataset/processed/eval_manifest_phrase_v2.json
% strategy: phrase
% split_seed: 42
% reporting_rows: 17942
% reporting_phrases: 9
% reporting_groups: 5
% max_length: 96
% git_commit: 0d29fb7
% generated_at: 2026-09-18T03:04:58.214778+00:00
%
\subsection{A single-metric safety gate certified a model that had abandoned an urgency class}
\label{sec:gate-failure}

Our acceptance gate was initially a single condition, CRITICAL recall $\geq 0.95$, on the reasoning that missing a critical presentation is the failure that matters. One configuration passed it with a CRITICAL recall of 0.9749 while achieving an URGENT recall of 0.0083: it assigned the CRITICAL label to 7,633 of 7,793 URGENT rows. It passed the safety gate \emph{because} it over-triaged, and the gate contained no term that could observe this.

The degeneracy is measurable rather than interpretive. A model that merges CRITICAL and URGENT earns a CRITICAL precision of $0.5349$ on this set by construction. The model scored $0.5337$, which is $0.0012$ \emph{below} its own set's merge-strategy precision. Within a thousandth, it did not approximate the merge strategy; it was the merge strategy.

Two further observations. First, the gate certified this model while rejecting a better one: a configuration with macro F1 0.7724 and all three classes alive failed the same gate on recall alone. Second, our own automated degeneracy check, written for exactly this failure, used the test $\mathrm{recall} \geq 0.9 \wedge \mathrm{precision} < 0.5$ and did not fire, because $0.5337$ sits above a hardcoded $0.5$ while the true merge-strategy precision for this set was $0.5349$. The lesson is not that the constant was too low: it is that the quantity is a property of the evaluation set and must be computed against it.



\subsection{The acceptance gate}

% GENERATED FILE - DO NOT EDIT BY HAND.
% Written by training/evaluate.py --writeup in ONE evaluation pass.
% Editing this file to change a reported number is fabrication;
% re-run the emitter instead.
%
% reported_model: model_v2d_freeze8_lr1e-5
% run_fingerprint: 3fdcb8c5d5f8f137
% manifest: dataset/processed/eval_manifest_phrase_v2.json
% strategy: phrase
% split_seed: 42
% reporting_rows: 17942
% reporting_phrases: 9
% reporting_groups: 5
% max_length: 96
% git_commit: ed02093
% generated_at: 2026-09-08T04:59:19.668469+00:00
%
\subsection{The acceptance gate, and where each threshold comes from}
\label{sec:gate}

A model is accepted only if all three conditions hold. They are not traded off against one another.

\paragraph{G1: CRITICAL recall $\geq 0.95$.} \emph{Inherited; source not verified.} This threshold predates the current record and carries only the note that missing a critical case is the failure that matters. It is consistent with the trauma field-triage convention of holding under-triage at or below 5\%, but we have not verified that attribution against a source and do not claim it. We keep the value because loosening a safety threshold on no evidence is worse than retaining an unsourced one.

\paragraph{G2: macro F1 $\geq 2/3$.} \emph{Derived exactly.} Macro F1 is the unweighted mean of three per-class F1 scores. If any one class is abandoned its F1 is zero, so macro F1 $\leq (1+1+0)/3 = 0.6667$ \emph{even when the other two classes are perfect}. Requiring macro F1 above $2/3$ is therefore a guarantee that no class has been abandoned, and it is the tightest such guarantee obtainable from a single scalar. This is a floor of non-degeneracy, not of quality.

\paragraph{G3: CRITICAL precision $\geq$ the degenerate ceiling plus a margin.} \emph{Floor derived and measured per evaluation set; margin not derived.} A model that merges CRITICAL and URGENT and labels the union CRITICAL earns, by construction, a CRITICAL precision of
\[ \frac{\mathrm{support(CRITICAL)}}{\mathrm{support(CRITICAL)} + \mathrm{support(URGENT)}} = 0.5349 \]
on this reporting set. That is what the degenerate strategy is paid for free, so any informative model must exceed it. We compute it against the set being scored rather than hardcoding a value, because it moves with the set's composition. \textbf{The margin above that floor (0.10) is not derived.} How far above provably-degenerate a deployable model must sit is a question about tolerable over-triage in a clinic, it is a clinical judgement that has not yet been made, and we report it as a placeholder rather than as a standard.


